\subsection{How much computation does the prefix certificate save?}
We developed the certificate after observing the largest-layer agreements and replayed it on all 55 training splits without using test labels. At the 10\% budget, all returned supports agreed both with the archived selections and with full evaluation under the same arithmetic and tie conventions. In 48 splits the first $h$ candidates sufficed. Even if two methods yield the same support, a bound tied at the cutoff may require evaluating extra candidates. Across the 55 splits the median fraction of all layer--column scores evaluated was 0.401. This count excludes the cost of finding neighbours and forming the layer weights.

The support implementation stores log binomial coefficients only for counts actually present in the witnesses and sums the positive and negative ballots without converting an entire dense mask to floating point. We compared three implementations: the original full scorer, the new scorer with all columns evaluated, and the same new scorer using the certificate. The last pair isolates pruning. At 16,384 columns and 192 fixed witnesses, the median times across six matrices were 44.73, 16.93 and 11.87\,ms, respectively. The median paired ratio for certified versus exhaustive evaluation within the new implementation was 0.72; at 64 columns it was 1.02. On fixed training witnesses, the certified minimum took a median 1.117 times the scoring time of the largest layer alone on Gisette and 1.194 times on ISOLET. All five supports agreed on each source, so certifying the full objective cost an additional 12\% and 19\%, respectively. When the run began with the original training array and included reconstructing the triplets, the paired ratios on Gisette and ISOLET were 0.94 and 0.97 against the new exhaustive scorer, and 0.78 and 0.79 against the original scorer. Thus most of the fixed-matrix saving did not carry through the cost of neighbour construction. These measurements stop when the support is returned and exclude downstream fitting.
